\documentclass[10pt,a4paper]{amsart}
\usepackage[margin=19mm]{geometry}
\usepackage{amsmath,amssymb,mathtools}
\usepackage[scr=rsfs,cal=euler]{mathalfa}
\usepackage{xfrac}
\usepackage[unicode,hidelinks,bookmarks=false]{hyperref}
\newcommand{\ie}{i.e.\ }
\newcommand{\cf}{cf.\ }
\newcommand{\resp}{respectively\ }
\newcommand{\Subsection}{\subsection}
\providecommand{\nobreakdash}{\nobreak\hbox{-}}
\setlength{\emergencystretch}{2em}
\multlinegap0pt
\usepackage{bm}
\usepackage{bbm}
\usepackage{dsfont}
\usepackage{xspace}
\usepackage{cancel}
\usepackage{mathtools}
\usepackage{amsthm}
\makeatletter
\@ifundefined{theorem}{\newtheorem{theorem}{Theorem}}{}
\makeatother
\title[Adjacency-degree algebras and spectral determination of grap]{Adjacency-degree algebras and spectral determination of graphs}
\author{Zhipeng Lu, Pengxiang Li}
\date{Source: arXiv:2607.21494v1 (CC BY 4.0)}
\begin{document}
\maketitle

\noindent\emph{Source and licence.} Adjacency-degree algebras and spectral determination of graphs. Version arXiv:2607.21494v1 (CC BY 4.0), distributed under the Creative Commons Attribution 4.0 International licence, as recorded in the arXiv licence metadata for this exact version and shown on the abstract page https://arxiv.org/abs/2607.21494v1. Full rights notice: licenses/arxiv-2607.21494-CC-BY-4.0.txt.\par\smallskip

\noindent\textit{Editorial scope.} Counted unit: the Theorem whose source label is \texttt{thm:tree-quotient-rigidity} in the article (source0.tex lines 862--865, source label \texttt{thm:tree-quotient-rigidity}), delivered together with the article's own complete original proof of it (source0.tex lines 867--899, one \texttt{proof} environment of 33 lines). No mathematics is written, repaired or re-derived: statement and proof are the article's own text line for line. Required context retained uncounted: none. Every definition, display and label of those lines is retained unchanged. Omitted from this package and not redistributed: 1--861, 866--866, 900--1409. Nothing inside a retained excerpt is omitted or altered: every retained body is delivered exactly as the article's lines are written, so no omission receipt of any kind accompanies this package. Citations: the retained excerpts carry no citation command at all, so no bibliography entry of the article is transcribed and no reference is redistributed. Reference adaptation: none was needed and none was made; every reference inside the delivered unit resolves inside it, so no unresolved reference or citation is produced. Printed slips kept literal: no printed word, symbol, hypothesis or number of the counted statement or of its proof was altered. Layout and class: the article's own class is \texttt{cas-sc} with packages including natbib, amsmath, amssymb, amsfonts, amsthm, mathtools, enumitem, booktabs; the wrapper is the lane's pinned \texttt{amsart} editorial wrapper at 10pt on A4, which restates the theorem-like environments the retained text needs and adds the article's own macro definitions that the retained text needs (none), with extra wrapper packages (bm, bbm, dsfont, xspace, cancel, mathtools). The delivered numbering is the unit's own. Assets: no retained span contains a figure environment, a graphics-inclusion command, a PDF-page inclusion, a table or a TikZ picture; no image file is copied into this package, the package declares no asset dependency and no image file is redistributed with it. Licence: version arXiv:2607.21494v1 of this article is distributed under the Creative Commons Attribution 4.0 International licence, as recorded in the arXiv licence metadata for this exact version and shown on the abstract page https://arxiv.org/abs/2607.21494v1. A full rights notice is delivered at licenses/arxiv-2607.21494-CC-BY-4.0.txt. Characters: every retained excerpt is pure ASCII, so the delivered engine, XeLaTeX with its default Latin Modern text font, reports no missing character.
\begin{theorem}[Orbit-quotient rigidity]
\label{thm:tree-quotient-rigidity}
A tree is determined up to isomorphism by its weighted orbit quotient.
\end{theorem}

\begin{proof}
Let \(Q=(s_i,b_{ij})\) be the weighted orbit quotient of a tree.  For an active
set \(S\subseteq\{1,\ldots,m\}\), put
\[
        d_i(S)=\sum_{j\in S} b_{ij}.
\]
Starting with \(S_0=\{1,\ldots,m\}\), delete all quotient vertices with
\(d_i(S_t)=1\), except when the remaining active quotient is one of the two
center-edge cases
\[
        S=\{c\},\quad s_c=2,\quad b_{cc}=1,
        \qquad\text{or}\qquad
        S=\{c,c'\},\quad b_{cc'}=b_{c'c}=1.
\]
This is exactly the image of simultaneous leaf stripping on the original tree,
because an active vertex in orbit \(O_i\) has active degree \(d_i(S_t)\).  Hence
the process determines the center type.  If \(i\) is deleted at stage \(t\), its
unique active neighbor is the unique \(p(i)\in S_t\) with \(b_{i,p(i)}=1\); this
is the parent orbit of \(i\).

Define rooted types from the leaves inward.  If \(i\) is a deleted leaf orbit,
set \(\tau_i=\bullet\).  After all children of \(i\) have been assigned types,
set
\[
        \tau_i=\bullet\Bigl(\{\tau_j^{\,b_{ij}}:p(j)=i\}\Bigr),
\]
where the notation records the multiset of rooted child types attached to one
vertex of orbit \(O_i\).  This recursion is forced by \(Q\).  At the terminal
core, either a single center type \(\tau_c\) remains, or two rooted types are
glued across a center edge, or two copies of the same rooted type are glued
across a symmetric center edge represented by a loop \(b_{cc}=1\).  Thus \(Q\)
constructs a unique tree up to isomorphism.
\end{proof}

\end{document}
